\documentclass[
    language=english,
    doctype=exam,
    institution=none,
]{../../omnilatex}

\RequirePackage{config/document-settings}

\begin{document}

\maketitle

\begin{question}[5]
    Define the term \textbf{algorithm} and explain the difference between
    \textit{time complexity} and \textit{space complexity}.  Give a brief
    example of each.
\end{question}
\begin{answer}[2cm]
\end{answer}

\begin{question}[5]
    Which of the following data structures is most appropriate for implementing
    a breadth-first search (BFS)?
    \begin{enumerate}[(a)]
        \item Stack
        \item Queue
        \item Heap
        \item Binary search tree
    \end{enumerate}
    Justify your answer in one sentence.
\end{question}
\begin{answer}[1.5cm]
\end{answer}

\begin{question}[10]
    Let $f(n) = 3n^2 + 5n + 7$.  Show that $f(n) = O(n^2)$ using the
    formal definition of Big-O notation.

    \textit{Hint:} Find constants $c > 0$ and $n_0 \geq 1$ such that
    $f(n) \leq c \cdot n^2$ for all $n \geq n_0$.
\end{question}
\begin{answer}[4cm]
\end{answer}

\begin{question}[20]
    Prove by mathematical induction that for all integers $n \geq 1$:
    \[
        \sum_{k=1}^{n} k^2 = \frac{n(n+1)(2n+1)}{6}.
    \]

    \textbf{(a)} State the base case and verify it.\hfill\textit{[5 marks]}

    \textbf{(b)} State the inductive hypothesis.\hfill\textit{[5 marks]}

    \textbf{(c)} Complete the inductive step.\hfill\textit{[10 marks]}
\end{question}
\begin{answer}[6cm]
\end{answer}

\begin{question}[15]
    Consider the following recurrence relation describing the running time of a
    divide-and-conquer algorithm:
    \[
        T(n) = 2\,T\!\left(\frac{n}{2}\right) + n, \qquad T(1) = 1.
    \]
    Using the Master Theorem, determine the asymptotic running time $T(n)$.
    Show all steps of your reasoning, including identification of the parameters
    $a$, $b$, and $f(n)$, and the comparison of $f(n)$ with $n^{\log_b a}$.
\end{question}
\begin{answer}[5cm]
\end{answer}

\begin{question}[15]
    \textbf{(a)} Describe the difference between a \textbf{min-heap} and a
    \textbf{max-heap}.\hfill\textit{[5 marks]}

    \textbf{(b)} Given the sequence $\langle 25, 17, 36, 2, 8, 19, 45 \rangle$,
    draw the min-heap that results after inserting these elements one by one
    (in the order given).\hfill\textit{[10 marks]}
\end{question}
\begin{answer}[5cm]
\end{answer}

\end{document}
